% Options for packages loaded elsewhere
\PassOptionsToPackage{unicode}{hyperref}
\PassOptionsToPackage{hyphens}{url}
%
\documentclass[
  11pt,
]{article}
\usepackage{amsmath,amssymb}
\usepackage{iftex}
\ifPDFTeX
  \usepackage[T1]{fontenc}
  \usepackage[utf8]{inputenc}
  \usepackage{textcomp} % provide euro and other symbols
\else % if luatex or xetex
  \usepackage{unicode-math} % this also loads fontspec
  \defaultfontfeatures{Scale=MatchLowercase}
  \defaultfontfeatures[\rmfamily]{Ligatures=TeX,Scale=1}
\fi
\usepackage{lmodern}
\ifPDFTeX\else
  % xetex/luatex font selection
\fi
% Use upquote if available, for straight quotes in verbatim environments
\IfFileExists{upquote.sty}{\usepackage{upquote}}{}
\IfFileExists{microtype.sty}{% use microtype if available
  \usepackage[]{microtype}
  \UseMicrotypeSet[protrusion]{basicmath} % disable protrusion for tt fonts
}{}
\makeatletter
\@ifundefined{KOMAClassName}{% if non-KOMA class
  \IfFileExists{parskip.sty}{%
    \usepackage{parskip}
  }{% else
    \setlength{\parindent}{0pt}
    \setlength{\parskip}{6pt plus 2pt minus 1pt}}
}{% if KOMA class
  \KOMAoptions{parskip=half}}
\makeatother
\usepackage{xcolor}
\usepackage[margin=2.5cm]{geometry}
\usepackage{graphicx}
\makeatletter
\newsavebox\pandoc@box
\newcommand*\pandocbounded[1]{% scales image to fit in text height/width
  \sbox\pandoc@box{#1}%
  \Gscale@div\@tempa{\textheight}{\dimexpr\ht\pandoc@box+\dp\pandoc@box\relax}%
  \Gscale@div\@tempb{\linewidth}{\wd\pandoc@box}%
  \ifdim\@tempb\p@<\@tempa\p@\let\@tempa\@tempb\fi% select the smaller of both
  \ifdim\@tempa\p@<\p@\scalebox{\@tempa}{\usebox\pandoc@box}%
  \else\usebox{\pandoc@box}%
  \fi%
}
% Set default figure placement to htbp
\def\fps@figure{htbp}
\makeatother
\setlength{\emergencystretch}{3em} % prevent overfull lines
\providecommand{\tightlist}{%
  \setlength{\itemsep}{0pt}\setlength{\parskip}{0pt}}
\setcounter{secnumdepth}{5}
\usepackage{float}
\usepackage{pdflscape}
\usepackage{array}
\usepackage{dcolumn}
\usepackage{booktabs}
\usepackage{placeins}
\usepackage{longtable}
\usepackage{multirow}
\usepackage{wrapfig}
\usepackage{colortbl}
\usepackage{tabu}
\usepackage{threeparttable}
\usepackage{threeparttablex}
\usepackage[normalem]{ulem}
\usepackage{makecell}
\usepackage{xcolor}
\usepackage{bookmark}
\IfFileExists{xurl.sty}{\usepackage{xurl}}{} % add URL line breaks if available
\urlstyle{same}
\hypersetup{
  pdftitle={Projekt Zaliczeniowy BTC},
  pdfauthor={Jakub Morawski, Jarosław Kwiatkowski},
  hidelinks,
  pdfcreator={LaTeX via pandoc}}

\title{Projekt Zaliczeniowy BTC}
\author{Jakub Morawski, Jarosław Kwiatkowski}
\date{2025-06-20}

\begin{document}
\maketitle

{
\setcounter{tocdepth}{2}
\tableofcontents
}
\section{Wstęp}\label{wstux119p}

\textbf{Kryptowaluty}, a zwłaszcza Bitcoin (BTC), stały się w ostatniej
dekadzie jednym z najbardziej fascynujących i dynamicznych zjawisk
finansowych. Ich wartość, charakteryzująca się wyjątkową zmiennością,
stanowi wyzwanie dla tradycyjnych modeli ekonomicznych. Co napędza
wahania cen najpopularniejszej kryptowaluty? Czy klucz leży w
fundamentalnych wskaźnikach technologicznych, makroekonomicznych
czynnikach globalnych, a może w irracjonalnych impulsach popkultury? To
pytania, na które niniejszy projekt próbuje odpowiedzieć poprzez
zastosowanie narzędzi ekonometrycznych.

\textbf{Celem projektu} jest zbudowanie i kompleksowa weryfikacja modelu
ekonometrycznego wyjaśniającego zmienność \textbf{średniej miesięcznej
ceny Bitcoina}. Wykorzystując szeregi czasowe, analizujemy okres
\texttt{1.08.2010\ -\ 1.05.2025} (dane miesięczne).

\textbf{Struktura sprawozdania} odzwierciedla pełny proces
modelowania:\\
1. Prezentację i diagnostykę danych\\
2. Transformacje zapewniające stacjonarność szeregów\\
3. Dobór zmiennych metodą \texttt{wybrana\_metoda}\\
4. Budowę modelu i jego weryfikację (testy RESET, Chowa, VIF,
autokorelacji, normalności reszt)\\
5. Interpretację parametrów i ocenę trafności modelu\\
6. Prognozę ex-post z obliczeniem błędów

Poniższa analiza stanowi praktyczną ilustrację wyzwań w modelowaniu
rynków krypto -- od problemu niestacjonarności przez wpływ szoków
medialnych po interpretację parametrów w środowisku ekstremalnej
zmienności. Wyniki nie tylko weryfikują hipotezy o czynnikach
wpływających na cenę BTC, ale także demonstrują siłę i ograniczenia
ekonometrii w analizie nowych klas aktywów.

\section{Prezentacja i opis wykorzystywanych
danych}\label{prezentacja-i-opis-wykorzystywanych-danych}

\subsection{Opis}\label{opis}

Jako zmienne objaśniające uwzględniliśmy zarówno czynniki
fundamentalne:\\
- \texttt{VolumeM} (miesięczny wolumen transakcji),\\
- \texttt{Difficulty} (trudność wydobycia bloków),\\
- \texttt{BTCCirculation} (liczbę BTC w obiegu),\\
- \texttt{USAvgKWHPrice} (średnią cenę energii elektrycznej w USA),\\
- \texttt{InflationRateUS} (stopę inflacji w USA),\\
jak i czynniki behawioralne:\\
- \texttt{GoogleTrends} (poziom wyszukiwań frazy ``bitcoin'' w
wyszukiwarce Google),\\
- \texttt{HasMuskMentioned} (liczbę tweetów Elona Muska dotyczących
kryptowalut).

\subsection{Źródła}\label{ux17aruxf3dux142a}

\textbf{Źródła danych} zostały starannie dobrane pod kątem wiarygodności
i kompletności:\\
- Dane technologiczne (cena, wolumen, trudność, podaż):
\href{https://bitinfocharts.com}{BitInfoCharts},
\href{https://www.coinwarz.com/mining/bitcoin/difficulty-chart}{CoinWarz}\\
- Dane ekonomiczne: Federal Reserve Bank of St.~Louis (FRED)\\
- Aktywność medialna: API X.com (dane Muska),
\href{https://trends.google.com}{Google Trends}

\subsection{Analiza podstawowych
statystyk}\label{analiza-podstawowych-statystyk}

Kompleksowa analiza statystyk opisowych stanowi niezbędny etap
przygotowania danych do modelowania ekonometrycznego. W przypadku
niestabilnych rynków kryptowalut, identyfikacja:

\begin{itemize}
\item
  charakteru rozkładów zmiennych
\item
  występowania wartości ekstremalnych
\item
  stopnia zmienności danych
\item
  potencjalnych problemów z normalnością rozkładu
\end{itemize}

pozwala na świadome podejmowanie decyzji dotyczących:

\begin{itemize}
\item
  Konieczności transformacji zmiennych (np. logarytmowanie)
\item
  Stosowania metod odpornych na autokorelację i heteroskedastyczność
\item
  Doboru odpowiednich testów statystycznych
\item
  Interpretacji parametrów modelu
\end{itemize}

Pominięcie tej fazy diagnostycznej może prowadzić do błędnej
specyfikacji modelu i nieprawidłowych wniosków ekonomicznych.

\subsubsection{Średnia i mediana}\label{ux15brednia-i-mediana}

Poniższa tabela przedstawia obliczone wartości średniej i mediany
wszystkich analizowanych zmiennych.

\begin{table}[!h]
\centering
\caption{\label{tab:unnamed-chunk-2}Miary tendencji centralnej}
\centering
\begin{tabular}[t]{lrr}
\toprule
Zmienna & Średnia & Mediana\\
\midrule
\cellcolor{gray!10}{Price} & \cellcolor{gray!10}{16412.5051} & \cellcolor{gray!10}{4231.4500}\\
VolumeM & 347.4345 & 2.1600\\
\cellcolor{gray!10}{USAvgKWHPrice} & \cellcolor{gray!10}{0.1429} & \cellcolor{gray!10}{0.1370}\\
Difficulty & 17477581170803.0625 & 1543670000000.0000\\
\cellcolor{gray!10}{BTCCirculation} & \cellcolor{gray!10}{15253093.2135} & \cellcolor{gray!10}{16741181.0000}\\
\addlinespace
InflationRateUS & 2.6062 & 2.1000\\
\cellcolor{gray!10}{HasMuskMentioned} & \cellcolor{gray!10}{0.1685} & \cellcolor{gray!10}{0.0000}\\
GoogleTrends & 1011829532.5843 & 410741000.0000\\
\bottomrule
\end{tabular}
\end{table}

\paragraph{\texorpdfstring{Obserwacje
\newline}{Obserwacje }}\label{obserwacje}

Analizując miary tendencji centralnej dla zmiennych modelu, można
zaobserwować następujące prawidłowości:

\begin{enumerate}
\def\labelenumi{\arabic{enumi}.}
\item
  \textbf{Cena Bitcoina (Price)} charakteryzuje się znaczną asymetrią -
  średnia wartość (16 412 USD) jest niemal czterokrotnie wyższa od
  mediany (4 231 USD), co wskazuje na występowanie okresów ekstremalnie
  wysokich wycen w analizowanym okresie.
\item
  \textbf{Wolumen transakcji (VolumeM)} prezentuje szczególnie
  interesujący rozkład, gdzie średnia (347.43) znacznie przewyższa
  medianę (2.16), co sugeruje istnienie nielicznych okresów bardzo
  intensywnej aktywności handlowej.
\item
  \textbf{Trudność wydobycia (Difficulty)} wykazuje kolosalną rozpiętość
  wartości (średnia: 17.5 bln, mediana: 1.5 bln), odzwierciedlając
  dynamiczny rozwój mocy obliczeniowej sieci Bitcoin.
\item
  \textbf{Zmienne behawioralne} prezentują szczególnie wysoką zmienność:

  \begin{itemize}
  \tightlist
  \item
    \textbf{Tweety Muska (HasMuskMentioned)} mają medianę równą 0 przy
    średniej 0.17, potwierdzając ich sporadyczny charakter
  \item
    \textbf{Wyszukiwania Google (GoogleTrends)} wykazują olbrzymią
    rozpiętość (średnia: 1 mld, mediana: 411 mln)
  \end{itemize}
\end{enumerate}

\paragraph{\texorpdfstring{Wnioski \newline}{Wnioski }}\label{wnioski}

Na podstawie analizy miar tendencji centralnej można sformułować
następujące konkluzje:

\begin{enumerate}
\def\labelenumi{\arabic{enumi}.}
\tightlist
\item
  \textbf{Charakter rynku kryptowalut}:

  \begin{itemize}
  \tightlist
  \item
    Ekstremalna zmienność cen i wolumenu transakcji wymaga zastosowania
    specjalnych metod analitycznych
  \item
    Tradycyjne miary statystyczne mogą nie w pełni oddawać specyfikę
    tego rynku
  \end{itemize}
\item
  \textbf{Implikacje dla modelowania}: \vspace{0.2cm}

  \begin{itemize}
  \tightlist
  \item
    Konieczność transformacji logarytmicznej zmiennych objaśniających
  \item
    Wskazane jest uwzględnienie efektów nieliniowych w specyfikacji
    modelu
  \item
    Należy rozważyć zastosowanie metod odpornych na wartości skrajne
  \end{itemize}
\item
  \textbf{Interpretacja ekonomiczna}: \vspace{0.2cm}

  \begin{itemize}
  \tightlist
  \item
    Czynniki fundamentalne wykazują większą stabilność niż behawioralne
  \item
    Mediana często lepiej reprezentuje typowe wartości niż średnia
  \item
    Wysoka zmienność czynników behawioralnych może mieć kluczowy wpływ
    na ceny
  \end{itemize}
\end{enumerate}

\vspace{0.3cm}

\textit{Dla pełniejszego obrazu zaleca się uzupełnienie tej analizy o:}

\begin{itemize}
\item Badanie rozkładów empirycznych poszczególnych zmiennych
\item Analizę współzależności między zmiennymi
\item Testy normalności rozkładów
\end{itemize}

\pagebreak

\subsubsection{Odchylenie standardowe i
wariancja}\label{odchylenie-standardowe-i-wariancja}

Poniższa tabela przedstawia obliczone wartości odchylenia standardowego
i wariancji rozkładu wszystkich analizowanych zmiennych.

\begin{table}[!h]
\centering
\caption{\label{tab:unnamed-chunk-3}Miary rozproszenia danych}
\centering
\begin{tabular}[t]{lrr}
\toprule
Zmienna & Odchylenie & Wariancja\\
\midrule
\cellcolor{gray!10}{Price} & \cellcolor{gray!10}{24372.5225} & \cellcolor{gray!10}{594019853.6838}\\
VolumeM & 3409.4042 & 11624036.8090\\
\cellcolor{gray!10}{USAvgKWHPrice} & \cellcolor{gray!10}{0.0155} & \cellcolor{gray!10}{0.0002}\\
Difficulty & 28508733777135.4258 & 812747901575582210644448888.0000\\
\cellcolor{gray!10}{BTCCirculation} & \cellcolor{gray!10}{4378359.6522} & \cellcolor{gray!10}{19170033244180.3398}\\
\addlinespace
InflationRateUS & 1.9753 & 3.9019\\
\cellcolor{gray!10}{HasMuskMentioned} & \cellcolor{gray!10}{0.5771} & \cellcolor{gray!10}{0.3330}\\
GoogleTrends & 1412356054.0194 & 1994749623325263104.0000\\
\bottomrule
\end{tabular}
\end{table}

\paragraph{\texorpdfstring{Obserwacje
\newline}{Obserwacje }}\label{obserwacje-1}

Analizując miary rozproszenia dla poszczególnych zmiennych, można
zaobserwować następujące prawidłowości:

\begin{enumerate}
\def\labelenumi{\arabic{enumi}.}
\item
  \textbf{Cena Bitcoina (Price)} charakteryzuje się ekstremalnie wysokim
  odchyleniem standardowym (24 372 USD) i wariancją (594 mln), co
  potwierdza ogromną zmienność cenową tego aktywa.
\item
  \textbf{Wolumen transakcji (VolumeM)} wykazuje kolosalne zróżnicowanie
  (odchylenie 3 409), znacznie przewyższające jego średnią wartość, co
  wskazuje na okresy skrajnie wysokiej i niskiej płynności.
\item
  \textbf{Trudność wydobycia (Difficulty)} prezentuje niewyobrażalnie
  wysokie wartości wariancji (\(\approx 8.13 \times 10^{23}\)), co
  wynika z wykładniczego wzrostu mocy obliczeniowej sieci Bitcoin.
\item
  \textbf{Zmienne makroekonomiczne}:

  \begin{itemize}
  \tightlist
  \item
    Cena energii (USAvgKWHPrice) ma stosunkowo niskie odchylenie
    (0.0155), typowe dla stabilnych wskaźników ekonomicznych
  \item
    Inflacja (InflationRateUS) wykazuje umiarkowaną zmienność
    (odchylenie 1.9753)
  \end{itemize}
\item
  \textbf{Zmienne behawioralne}:

  \begin{itemize}
  \tightlist
  \item
    Zmieenna określająca liczbę Tweetów Elona Muska (HasMuskMentioned)
    ma odchylenie (0.5771) ponad 3-krotnie wyższe od średniej
  \item
    Wyszukiwania Google (GoogleTrends) charakteryzują się gigantyczną
    wariancją \((\approx 1.99 \times 10^{18})\)
  \end{itemize}
\end{enumerate}

\paragraph{Wnioski}\label{wnioski-1}

\begin{enumerate}
\def\labelenumi{\arabic{enumi}.}
\tightlist
\item
  \textbf{Charakterystyka ryzyka}: \vspace{0.2cm}

  \begin{itemize}
  \tightlist
  \item
    Ekstremalne wartości wariancji dla ceny BTC i wolumenu potwierdzają
    wysokie ryzyko inwestycyjne
  \item
    Niestabilność czynników behawioralnych może dominować nad
    fundamentami
  \end{itemize}
\item
  \textbf{Implikacje dla modelu}: \vspace{0.2cm}

  \begin{itemize}
  \tightlist
  \item
    Konieczność uwzględnienia heteroskedastyczności w specyfikacji
    modelu
  \item
    Wskazane jest zastosowanie transformacji logarytmicznej dla
    zmiennych o wysokiej wariancji
  \item
    Należy rozważyć użycie modeli GARCH dla zmiennej zależnej (Price)
  \end{itemize}
\item
  \textbf{Interpretacja ekonomiczna}: \vspace{0.2cm}

  \begin{itemize}
  \tightlist
  \item
    Stabilność zmiennych makro (energia, inflacja) kontrastuje z
    niestabilnością rynku kryptowalutowego
  \item
    Czynniki technologiczne (Difficulty) wykazują nieliniowy charakter
    wzrostu
  \item
    Wpływ czynników behawioralnych może mieć charakter impulsowy
  \end{itemize}
\end{enumerate}

\vspace{0.3cm}

\textit{Uwaga: Ze względu na ekstremalne wartości niektórych miar rozproszenia, zaleca się:}

\begin{itemize}
\item Weryfikację założeń o normalności rozkładów
\item Analizę korelacji między zmiennymi objaśniającymi
\item Stosowanie metod odpornych na wartości odstające
\end{itemize}

\pagebreak

\subsubsection{Kurtoza}\label{kurtoza}

Poniższa tabela przedstawia obliczone wartości kurtozy rozkładu
wszystkich analizowanych zmiennych.

\begin{table}[!h]
\centering
\caption{\label{tab:unnamed-chunk-4}Spłaszczenie rozkładu (kurtoza)}
\centering
\begin{tabular}[t]{lr}
\toprule
Zmienna & Kurtoza\\
\midrule
\cellcolor{gray!10}{Price} & \cellcolor{gray!10}{2.152}\\
VolumeM & 144.686\\
\cellcolor{gray!10}{USAvgKWHPrice} & \cellcolor{gray!10}{0.158}\\
Difficulty & 2.899\\
\cellcolor{gray!10}{BTCCirculation} & \cellcolor{gray!10}{-0.081}\\
\addlinespace
InflationRateUS & 1.986\\
\cellcolor{gray!10}{HasMuskMentioned} & \cellcolor{gray!10}{30.485}\\
GoogleTrends & 8.625\\
\bottomrule
\end{tabular}
\end{table}

\paragraph{\texorpdfstring{Obserwacje
\newline}{Obserwacje }}\label{obserwacje-2}

Analizując wartości kurtozy dla poszczególnych zmiennych, można
zaobserwować następujące prawidłowości:

\begin{enumerate}
\def\labelenumi{\arabic{enumi}.}
\item
  \textbf{Cena Bitcoina (Price)} wykazuje umiarkowaną leptokurtyczność
  (kurtoza = 2.152), co wskazuje na nieco bardziej strome i
  ``ciężkoogonowe'' rozkłady niż rozkład normalny.
\item
  \textbf{Wolumen transakcji (VolumeM)} charakteryzuje się ekstremalną
  leptokurtycznością (144.686), co sugeruje występowanie znacznej
  koncentracji wartości wokół średniej przy jednoczesnym występowaniu
  skrajnych wartości odstających.
\item
  \textbf{Cena energii (USAvgKWHPrice)} jest prawie mezokurtyczna
  (0.158), co oznacza rozkład zbliżony do normalnego, typowy dla
  stabilnych wskaźników ekonomicznych.
\item
  \textbf{Trudność wydobycia (Difficulty)} wykazuje wyraźną
  leptokurtyczność (2.899), co odzwierciedla okresowe skoki wartości
  związane z dostosowywaniem trudności wydobycia.
\item
  \textbf{Liczba BTC w obiegu (BTCCirculation)} jest jedyną zmienną o
  ujemnej kurtozie (-0.081), co wskazuje na bardziej płaski rozkład niż
  normalny.
\item
  \textbf{Zmienne behawioralne}:

  \begin{itemize}
  \tightlist
  \item
    Tweety Muska (HasMuskMentioned) wykazują skrajną leptokurtyczność
    (30.485)
  \item
    Wyszukiwania Google (GoogleTrends) mają wysoką leptokurtyczność
    (8.625)
  \end{itemize}
\end{enumerate}

\paragraph{Wnioski}\label{wnioski-2}

\begin{enumerate}
\def\labelenumi{\arabic{enumi}.}
\tightlist
\item
  \textbf{Interpretacja kształtu rozkładów}: \vspace{0.2cm}

  \begin{itemize}
  \tightlist
  \item
    Większość zmiennych wykazuje leptokurtyczność, co sugeruje częstsze
    występowanie ekstremów niż w rozkładzie normalnym
  \item
    Ekstremalne wartości kurtozy dla wolumenu i tweetów Muska wskazują
    na impulsowy charakter tych zmiennych
  \end{itemize}
\item
  \textbf{Implikacje dla modelowania}: \vspace{0.2cm}

  \begin{itemize}
  \tightlist
  \item
    Wysoka kurtoza wymaga zastosowania metod odpornych na wartości
    odstające
  \item
    Konieczność transformacji zmiennych o skrajnej leptokurtyczności
  \item
    Należy rozważyć modele uwzględniające ``grube ogony'' rozkładów
  \end{itemize}
\item
  \textbf{Znaczenie ekonomiczne}: \vspace{0.2cm}

  \begin{itemize}
  \tightlist
  \item
    Leptokurtyczne rozkłady potwierdzają wysoką niestabilność rynku
    kryptowalut
  \item
    Stabilność cen energii kontrastuje z niestabilnością zmiennych
    rynkowych
  \item
    Impulsowy charakter czynników behawioralnych może dominować w
    krótkim okresie
  \end{itemize}
\end{enumerate}

\vspace{0.3cm}

\textit{Uwaga: W przypadku zmiennych o wysokiej kurtozie zaleca się:}

\begin{itemize}
\item Weryfikację normalności rozkładu testami statystycznymi
\item Rozważenie transformacji Boxa-Coxa
\item Analizę wartości odstających
\end{itemize}

\pagebreak

\subsubsection{Skośność}\label{skoux15bnoux15bux107}

\paragraph{\texorpdfstring{Obliczenie
\newline}{Obliczenie }}\label{obliczenie}

Poniższa tabela przedstawia obliczone wartości współczynnika skośności
rozkładu wszystkich analizowanych zmiennych.

\begin{table}[!h]
\centering
\caption{\label{tab:unnamed-chunk-5}Współczynnik skośności rozkładu}
\centering
\begin{tabular}[t]{lr}
\toprule
Zmienna & Skośność\\
\midrule
\cellcolor{gray!10}{Price} & \cellcolor{gray!10}{1.706}\\
VolumeM & 11.793\\
\cellcolor{gray!10}{USAvgKWHPrice} & \cellcolor{gray!10}{1.243}\\
Difficulty & 1.934\\
\cellcolor{gray!10}{BTCCirculation} & \cellcolor{gray!10}{-0.991}\\
\addlinespace
InflationRateUS & 1.518\\
\cellcolor{gray!10}{HasMuskMentioned} & \cellcolor{gray!10}{4.892}\\
GoogleTrends & 2.627\\
\bottomrule
\end{tabular}
\end{table}

\paragraph{Obserwacje}\label{obserwacje-3}

Analiza współczynników skośności dla badanych zmiennych ujawnia
następujące prawidłowości:

\begin{enumerate}
\def\labelenumi{\arabic{enumi}.}
\item
  \textbf{Cena Bitcoina (Price)} wykazuje silną prawostronną asymetrię
  (skośność = 1.706), co potwierdza dominację okresów wzrostowych z
  rzadkimi, ale gwałtownymi spadkami.
\item
  \textbf{Wolumen transakcji (VolumeM)} charakteryzuje się ekstremalną
  skośnością dodatnią (11.793), wskazując na występowanie nielicznych
  okresów anomalnie wysokiej aktywności handlowej.
\item
  \textbf{Cena energii (USAvgKWHPrice)} prezentuje umiarkowaną skośność
  dodatnią (1.243), typową dla stabilnych wskaźników ekonomicznych z
  okresowymi wzrostami.
\item
  \textbf{Trudność wydobycia (Difficulty)} również wykazuje silną
  asymetrię prawostronną (1.934), odzwierciedlającą stopniowy, ale
  nieliniowy wzrost mocy obliczeniowej sieci.
\item
  \textbf{Liczba BTC w obiegu (BTCCirculation)} jest jedyną zmienną o
  ujemnej skośności (-0.991), co sugeruje okresowe zmniejszanie się
  podaży w analizowanym okresie.
\item
  \textbf{Zmienne behawioralne}:

  \begin{itemize}
  \tightlist
  \item
    Tweety Muska (HasMuskMentioned) wykazują bardzo silną skośność
    dodatnią (4.892)
  \item
    Wyszukiwania Google (GoogleTrends) mają wyraźną asymetrię
    prawostronną (2.627)
  \end{itemize}
\end{enumerate}

\paragraph{Wnioski}\label{wnioski-3}

\begin{enumerate}
\def\labelenumi{\arabic{enumi}.}
\tightlist
\item
  \textbf{Interpretacja ekonomiczna}: \vspace{0.2cm}

  \begin{itemize}
  \tightlist
  \item
    Dominacja dodatniej skośności potwierdza charakterystyczną dla rynku
    kryptowalunt nierównowagę między okresami hossy i bessy
  \item
    Ujemna skośność podaży BTC może wynikać z mechanizmu halvingu
    ograniczającego emisję nowych monet
  \end{itemize}
\item
  \textbf{Implikacje metodologiczne}: \vspace{0.2cm}

  \begin{itemize}
  \tightlist
  \item
    Wymóg transformacji zmiennych o wysokiej skośności przed
    modelowaniem
  \item
    Konieczność uwzględnienia nieliniowości w specyfikacji modelu
  \item
    Zalecenie stosowania testów statystycznych odpornych na asymetrię
  \end{itemize}
\item
  \textbf{Znaczenie dla prognozowania}: \vspace{0.2cm}

  \begin{itemize}
  \tightlist
  \item
    Asymetria rozkładów sugeruje różne właściwości statystyczne w
    okresach wzrostów i spadków
  \item
    Wysoka skośność czynników behawioralnych wskazuje na ich impulsowy
    charakter
  \item
    Należy rozważyć osobne modelowanie dla różnych faz rynkowych
  \end{itemize}
\end{enumerate}

\vspace{0.3cm}

\textit{Rekomendacje analityczne:}

\begin{itemize}
\item Dla zmiennych o skośności >1 zastosować transformację logarytmiczną
\item Przeprowadzić analizę podprób dla okresów o różnej dynamice
\item Uwzględnić w modelu efekty nieliniowe i interakcje
\end{itemize}

\pagebreak

\subsubsection{Współczynnik zmienności
(CV)}\label{wspuxf3ux142czynnik-zmiennoux15bci-cv}

Zgodnie z powszechnie przyjętą skalą klasyfikacji zmienności:

\begin{itemize}
\tightlist
\item
  CV \textless{} 15\% - zmienność niska
\item
  15\% \textless{} CV \textless{} 30\% - zmienność umiarkowana
\item
  30\% \textless{} CV \textless{} 100\% - zmienność wysoka
\item
  CV \textgreater{} 100\% - zmienność ekstremalna
\end{itemize}

\paragraph{\texorpdfstring{Obliczenie
\newline}{Obliczenie }}\label{obliczenie-1}

Poniższa tabela przedstawia obliczone wartości współczynnika zmienności
{[}\%{]} dla wszystkich analizowanych zmiennych.

\begin{table}[!h]
\centering
\caption{\label{tab:cv_table}Współczynnik zmienności (CV) [w procentach]}
\centering
\begin{tabular}[t]{lr}
\toprule
Zmienna & CV\\
\midrule
\cellcolor{gray!10}{Price} & \cellcolor{gray!10}{148.50}\\
VolumeM & 981.31\\
\cellcolor{gray!10}{USAvgKWHPrice} & \cellcolor{gray!10}{10.86}\\
Difficulty & 163.12\\
\cellcolor{gray!10}{BTCCirculation} & \cellcolor{gray!10}{28.70}\\
\addlinespace
InflationRateUS & 75.79\\
\cellcolor{gray!10}{HasMuskMentioned} & \cellcolor{gray!10}{342.40}\\
GoogleTrends & 139.58\\
\bottomrule
\end{tabular}
\end{table}

\paragraph{Obserwacje}\label{obserwacje-4}

Analiza współczynników zmienności (CV) dla badanych zmiennych ujawnia
następujące prawidłowości:

\begin{enumerate}
\def\labelenumi{\arabic{enumi}.}
\item
  \textbf{Cena Bitcoina (Price)} charakteryzuje się bardzo wysoką
  zmiennością (CV=148.5\%), co jest typowe dla aktywów kryptowalutowych
  i odzwierciedla ich spekulacyjny charakter.
\item
  \textbf{Wolumen transakcji (VolumeM)} wykazuje ekstremalną zmienność
  (CV=981.31\%), najwyższą ze wszystkich analizowanych zmiennych,
  wskazując na ogromne różnice w aktywności handlowej między okresami.
\item
  \textbf{Cena energii (USAvgKWHPrice)} jest najbardziej stabilną
  zmienną (CV=10.86\%), co potwierdza jej charakter jako wskaźnika
  makroekonomicznego.
\item
  \textbf{Trudność wydobycia (Difficulty)} również cechuje wysoka
  zmienność (CV=163.12\%), odzwierciedlająca dynamiczne dostosowania w
  sieci Bitcoin.
\item
  \textbf{Zmienne behawioralne}:

  \begin{itemize}
  \tightlist
  \item
    Tweety Muska (HasMuskMentioned) wykazują skrajną zmienność
    (CV=342.4\%)
  \item
    Wyszukiwania Google (GoogleTrends) mają bardzo wysoką zmienność
    (CV=139.58\%)
  \end{itemize}
\item
  \textbf{Podaż BTC (BTCCirculation)} jest względnie stabilna
  (CV=28.7\%), co wynika z przewidywalnego mechanizmu emisji.
\end{enumerate}

\paragraph{Wnioski}\label{wnioski-4}

\begin{enumerate}
\def\labelenumi{\arabic{enumi}.}
\tightlist
\item
  \textbf{Interpretacja ekonomiczna}: \vspace{0.2cm}

  \begin{itemize}
  \tightlist
  \item
    Ekstremalna zmienność rynku kryptowalutowego kontrastuje ze
    stabilnością tradycyjnych wskaźników ekonomicznych
  \item
    Czynniki behawioralne są bardziej zmienne niż fundamentalne
  \item
    Stabilność podaży BTC potwierdza przewidywalność protokołu Bitcoin
  \end{itemize}
\item
  \textbf{Implikacje metodologiczne}: \vspace{0.2cm}

  \begin{itemize}
  \tightlist
  \item
    Konieczność zastosowania modeli uwzględniających
    heteroskedastyczność
  \item
    Wskazane jest użycie ważonych estymatorów dla zmiennych o różnej
    zmienności
  \item
    Należy rozważyć osobne modelowanie okresów wysokiej i niskiej
    zmienności
  \end{itemize}
\item
  \textbf{Znaczenie dla zarządzania ryzykiem}: \vspace{0.2cm}

  \begin{itemize}
  \tightlist
  \item
    Wysoka zmienność wymaga odpowiednich strategii hedgingowych
  \item
    Czynniki o CV\textgreater100\% mogą dominować w krótkoterminowych
    fluktuacjach cen
  \item
    Stabilne zmienne lepiej wyjaśniają trendy długoterminowe
  \end{itemize}
\end{enumerate}

\vspace{0.3cm}

\textit{Rekomendacje analityczne:}

\begin{itemize}
\item Dla zmiennych o CV>100\% zastosować modele GARCH/ARCH
\item Przeprowadzić analizę przedziałową dla różnych poziomów zmienności
\item Uwzględnić interakcje między zmiennymi o różnej zmienności
\end{itemize}

\pagebreak

\subsubsection{Minimum i maksimum}\label{minimum-i-maksimum}

\paragraph{\texorpdfstring{Obliczenie
\newline}{Obliczenie }}\label{obliczenie-2}

Poniższa tabela przedstawia największe i największe wartości wszystkich
analizowanych zmiennych.

\begin{table}[!h]
\centering
\caption{\label{tab:unnamed-chunk-6}Wartości minimalne i maksymalne}
\centering
\begin{tabular}[t]{lrr}
\toprule
Zmienna & Min & Max\\
\midrule
\cellcolor{gray!10}{Price} & \cellcolor{gray!10}{0.10} & \cellcolor{gray!10}{103293.40}\\
VolumeM & 0.06 & 43700.00\\
\cellcolor{gray!10}{USAvgKWHPrice} & \cellcolor{gray!10}{0.12} & \cellcolor{gray!10}{0.18}\\
Difficulty & 112.38 & 120975000000000.00\\
\cellcolor{gray!10}{BTCCirculation} & \cellcolor{gray!10}{3561446.00} & \cellcolor{gray!10}{19857514.00}\\
\addlinespace
InflationRateUS & -0.20 & 9.10\\
\cellcolor{gray!10}{HasMuskMentioned} & \cellcolor{gray!10}{0.00} & \cellcolor{gray!10}{5.00}\\
GoogleTrends & 497700.00 & 9351710000.00\\
\bottomrule
\end{tabular}
\end{table}

\paragraph{Obserwacje}\label{obserwacje-5}

Analiza wartości ekstremalnych dla badanych zmiennych ujawnia
następujące prawidłowości:

\begin{enumerate}
\def\labelenumi{\arabic{enumi}.}
\item
  \textbf{Cena Bitcoina (Price)} wykazuje niezwykły zakres wahań od
  \$0.10 do \$103,293.40, dokumentując zarówno początkową fazę rozwoju
  (2010 r.), jak i okresy szczytowych wycen.
\item
  \textbf{Wolumen transakcji (VolumeM)} prezentuje ekstremalną
  rozpiętość od 0.06 do 43,700, odzwierciedlającą zarówno okresy
  minimalnej aktywności, jak i gigantyczne wolumeny handlowe.
\item
  \textbf{Cena energii (USAvgKWHPrice)} charakteryzuje się wąskim
  zakresem (0.12-0.18 USD/kWh), typowym dla stabilnych wskaźników
  infrastrukturalnych.
\item
  \textbf{Trudność wydobycia (Difficulty)} pokazuje astronomiczny wzrost
  od 112.38 do 120.975 bln, ilustrujący ewolucję sieci Bitcoin od
  początków do obecnej skali.
\item
  \textbf{Zmienne behawioralne}:

  \begin{itemize}
  \tightlist
  \item
    Tweety Muska (HasMuskMentioned) obejmują zakres od 0 do 5 wpisów
    miesięcznie
  \item
    Wyszukiwania Google (GoogleTrends) zmieniają się od 497,700 do 9.35
    mld
  \end{itemize}
\item
  \textbf{Podaż BTC (BTCCirculation)} wzrosła z 3.56 mln do 19.86 mln
  BTC, zgodnie z przewidywalnym harmonogramem emisji.
\end{enumerate}

\paragraph{Wnioski}\label{wnioski-5}

\begin{enumerate}
\def\labelenumi{\arabic{enumi}.}
\tightlist
\item
  \textbf{Interpretacja historyczna}: \vspace{0.2cm}

  \begin{itemize}
  \tightlist
  \item
    Ekstremalne wartości minimalne dokumentują początkową fazę rozwoju
    rynku krypto
  \item
    Wartości maksymalne odzwierciedlają okresy szczytowego
    zainteresowania i hossy
  \end{itemize}
\item
  \textbf{Implikacje metodologiczne}: \vspace{0.2cm}

  \begin{itemize}
  \tightlist
  \item
    Konieczność uwzględnienia nieliniowości i efektów progowych
  \item
    Wskazane jest zastosowanie transformacji zmniejszających wpływ
    wartości skrajnych
  \item
    Należy rozważyć analizę przedziałową dla skrajnie różnych warunków
    rynkowych
  \end{itemize}
\item
  \textbf{Znaczenie dla modelowania}: \vspace{0.2cm}

  \begin{itemize}
  \tightlist
  \item
    Wartości ekstremalne mogą wymagać specyficznego podejścia do
    estymacji
  \item
    Należy testować stabilność parametrów dla różnych przedziałów
    wartości
  \item
    Wskazane jest uwzględnienie zmiennych dummy dla okresów
    ekstremalnych
  \end{itemize}
\end{enumerate}

\vspace{0.3cm}

\textit{Rekomendacje analityczne:}

\begin{itemize}
\item Przeprowadzić analizę outlierów i rozważyć ich ewentualną eliminację
\item Zastosować transformacje logarytmiczne dla zmiennych o dużych rozpiętościach
\item Rozważyć podział próby na podokresy o zbliżonych charakterystykach
\end{itemize}

\pagebreak

\subsubsection{Rozstęp międzykwartylowy
(IQR)}\label{rozstux119p-miux119dzykwartylowy-iqr}

\paragraph{\texorpdfstring{Obliczenie
\newline}{Obliczenie }}\label{obliczenie-3}

Poniższa tabela przedstawia obliczone wartości rozstępów
międzykwartylowych wszystkich analizowanych zmiennych.

\begin{table}[!h]
\centering
\caption{\label{tab:unnamed-chunk-7}Rozstęp międzykwartylowy (IQR)}
\centering
\begin{tabular}[t]{lr}
\toprule
Zmienna & IQR\\
\midrule
\cellcolor{gray!10}{Price} & \cellcolor{gray!10}{25032.2250}\\
VolumeM & 3.8200\\
\cellcolor{gray!10}{USAvgKWHPrice} & \cellcolor{gray!10}{0.0090}\\
Difficulty & 21562639210073.0000\\
\cellcolor{gray!10}{BTCCirculation} & \cellcolor{gray!10}{6183664.7500}\\
\addlinespace
InflationRateUS & 1.5750\\
\cellcolor{gray!10}{HasMuskMentioned} & \cellcolor{gray!10}{0.0000}\\
GoogleTrends & 1048012500.0000\\
\bottomrule
\end{tabular}
\end{table}

\paragraph{Obserwacje}\label{obserwacje-6}

Analiza rozstępów międzykwartylowych (IQR) dla badanych zmiennych
ujawnia następujące prawidłowości:

\begin{enumerate}
\def\labelenumi{\arabic{enumi}.}
\item
  \textbf{Cena Bitcoina (Price)} charakteryzuje się bardzo szerokim IQR
  (25,032.23), co potwierdza ogromną zmienność cen nawet w środkowych
  50\% rozkładu.
\item
  \textbf{Wolumen transakcji (VolumeM)} wykazuje stosunkowo wąski IQR
  (3.82) w porównaniu do ekstremalnych wartości minimalnych i
  maksymalnych, co sugeruje koncentrację większości obserwacji w wąskim
  zakresie.
\item
  \textbf{Cena energii (USAvgKWHPrice)} ma wyjątkowo wąski IQR (0.009),
  co potwierdza jej stabilność jako zmiennej makroekonomicznej.
\item
  \textbf{Trudność wydobycia (Difficulty)} prezentuje astronomiczny IQR
  (21.56 bln), odzwierciedlający szybki rozwój mocy obliczeniowej sieci.
\item
  \textbf{Zmienne behawioralne}:

  \begin{itemize}
  \tightlist
  \item
    Tweety Muska (HasMuskMentioned) mają IQR=0, co wskazuje na brak
    zmienności w środkowych 50\% danych
  \item
    Wyszukiwania Google (GoogleTrends) wykazują bardzo szeroki IQR (1.05
    mld)
  \end{itemize}
\item
  \textbf{Podaż BTC (BTCCirculation)} ma umiarkowany IQR (6.18 mln),
  zgodny z przewidywalnym harmonogramem emisji.
\end{enumerate}

\paragraph{Wnioski}\label{wnioski-6}

\begin{enumerate}
\def\labelenumi{\arabic{enumi}.}
\tightlist
\item
  \textbf{Interpretacja statystyczna}: \vspace{0.2cm}

  \begin{itemize}
  \tightlist
  \item
    Wysokie wartości IQR dla ceny i trudności wydobycia potwierdzają ich
    dynamiczny charakter
  \item
    Zerowy IQR dla tweetów Muska wskazuje na ich dyskretny, impulsowy
    charakter
  \item
    Niskie IQR cen energii potwierdza jej stabilność
  \end{itemize}
\item
  \textbf{Implikacje metodologiczne}: \vspace{0.2cm}

  \begin{itemize}
  \tightlist
  \item
    Zmienne o wysokim IQR wymagają specjalnych metod analitycznych
  \item
    Należy rozważyć analizę oddzielnie dla każdego kwartyla
  \item
    Wskazane jest zastosowanie metod odpornych na ekstremy
  \end{itemize}
\item
  \textbf{Znaczenie praktyczne}: \vspace{0.2cm}

  \begin{itemize}
  \tightlist
  \item
    Szerokie IQR zmiennych rynkowych wymaga elastycznych strategii
    inwestycyjnych
  \item
    Wąskie IQR zmiennych makro sugeruje ich ograniczoną przydatność do
    krótkoterminowego prognozowania
  \item
    Zerowy IQR tweetów Muska potwierdza ich nieregularny wpływ na rynek
  \end{itemize}
\end{enumerate}

\vspace{0.3cm}

\textit{Rekomendacje analityczne:}

\begin{itemize}
\item Dla zmiennych o IQR=0 rozważyć modelowanie dyskretne
\item Dla szerokich IQR zastosować transformacje normalizujące
\item Przeprowadzić analizę percentylową dla pełniejszego obrazu rozkładów
\end{itemize}

\pagebreak

\subsubsection{Medianowe odchylenie bezwględne
(MAD)}\label{medianowe-odchylenie-bezwglux119dne-mad}

\paragraph{\texorpdfstring{Obliczenie
\newline}{Obliczenie }}\label{obliczenie-4}

Poniższa tabela przedstawia obliczone wartości medianowego odchylenia
bezwzględnego wszystkich analizowanych zmiennych.

\begin{table}[!h]
\centering
\caption{\label{tab:unnamed-chunk-8}Medianowe odchylenie bezwzględne (MAD)}
\centering
\begin{tabular}[t]{lr}
\toprule
Zmienna & MAD\\
\midrule
\cellcolor{gray!10}{Price} & \cellcolor{gray!10}{6266.1348}\\
VolumeM & 1.8481\\
\cellcolor{gray!10}{USAvgKWHPrice} & \cellcolor{gray!10}{0.0059}\\
Difficulty & 2288645140244.8887\\
\cellcolor{gray!10}{BTCCirculation} & \cellcolor{gray!10}{3655341.4044}\\
\addlinespace
InflationRateUS & 1.1861\\
\cellcolor{gray!10}{HasMuskMentioned} & \cellcolor{gray!10}{0.0000}\\
GoogleTrends & 537811667.4000\\
\bottomrule
\end{tabular}
\end{table}

\paragraph{Obserwacje}\label{obserwacje-7}

Analiza medianowego odchylenia bezwzględnego (MAD) dla badanych
zmiennych ujawnia następujące prawidłowości:

\begin{enumerate}
\def\labelenumi{\arabic{enumi}.}
\item
  \textbf{Cena Bitcoina (Price)} wykazuje wysoką odporną miarę
  zmienności (MAD=6,266.13), co potwierdza znaczną rozpiętość cen nawet
  po wyeliminowaniu wpływu wartości skrajnych.
\item
  \textbf{Wolumen transakcji (VolumeM)} charakteryzuje się stosunkowo
  niskim MAD (1.85), co wskazuje na koncentrację większości obserwacji w
  pobliżu mediany przy jednoczesnych ekstremalnych wartościach
  odstających.
\item
  \textbf{Cena energii (USAvgKWHPrice)} prezentuje minimalne MAD
  (0.0059), potwierdzając jej stabilność jako wskaźnika
  makroekonomicznego.
\item
  \textbf{Trudność wydobycia (Difficulty)} ma ekstremalnie wysokie MAD
  (2.29 bln), odzwierciedlające ogromną skalę zmian w mocy obliczeniowej
  sieci.
\item
  \textbf{Zmienne behawioralne}:

  \begin{itemize}
  \tightlist
  \item
    Tweety Muska (HasMuskMentioned) mają MAD=0, co wskazuje na brak
    zmienności wokół mediany
  \item
    Wyszukiwania Google (GoogleTrends) wykazują bardzo wysokie MAD
    (537.8 mln)
  \end{itemize}
\item
  \textbf{Podaż BTC (BTCCirculation)} prezentuje umiarkowane MAD (3.66
  mln), zgodne z przewidywalnym charakterem emisji nowych monet.
\end{enumerate}

\paragraph{Wnioski}\label{wnioski-7}

\begin{enumerate}
\def\labelenumi{\arabic{enumi}.}
\tightlist
\item
  \textbf{Interpretacja statystyczna}: \vspace{0.2cm}

  \begin{itemize}
  \tightlist
  \item
    Wysokie wartości MAD dla ceny i trudności wydobycia potwierdzają ich
    inherentną zmienność
  \item
    Zerowe MAD tweetów Muska dokumentuje ich dyskretny charakter
  \item
    Niskie MAD cen energii wzmacnia jej przydatność jako stabilnej
    zmiennej objaśniającej
  \end{itemize}
\item
  \textbf{Zalety MAD jako miary}: \vspace{0.2cm}

  \begin{itemize}
  \tightlist
  \item
    Odporność na wartości odstające w porównaniu do odchylenia
    standardowego
  \item
    Lepsza reprezentacja typowej zmienności dla rozkładów o grubych
    ogonach
  \item
    Spójność z medianą jako miarą centralną
  \end{itemize}
\item
  \textbf{Implikacje dla modelowania}: \vspace{0.2cm}

  \begin{itemize}
  \tightlist
  \item
    Zmienne o wysokim MAD wymagają specjalnych technik estymacji
  \item
    Należy rozważyć ważenie obserwacji odwrotnie proporcjonalnie do MAD
  \item
    Wskazane jest testowanie modeli na podzbiorach danych wyznaczonych
    przez MAD
  \end{itemize}
\end{enumerate}

\vspace{0.3cm}

\textit{Rekomendacje analityczne:}

\begin{itemize}
\item Dla zmiennych o MAD=0 zastosować modele dyskretne lub dychotomiczne
\item Dla wysokich MAD wykorzystać metody statystyki robustowej
\item Przeprowadzić analizę wrażliwości dla różnych poziomów obcięcia rozkładów
\end{itemize}

\pagebreak

\subsubsection{Współczynnik Giniego}\label{wspuxf3ux142czynnik-giniego}

Zgodnie z powszechnie przyjętą skalą klasyfikacji wartości współczynnika
Giniego:

\begin{itemize}
\tightlist
\item
  0 - 0.3: niska koncentracja
\item
  0.3 - 0.6: umiarkowana koncentracja
\item
  0.6 - 0.8: wysoka koncentracja
\item
  0.8 - 1: ekstremalna koncentracja
\end{itemize}

\paragraph{\texorpdfstring{Obliczenie
\newline}{Obliczenie }}\label{obliczenie-5}

Poniższa tabela przedstawia obliczone wartości współczynnika
współczynnika koncentracji Giniego wszystkich analizowanych zmiennych.

\begin{table}[!h]
\centering
\caption{\label{tab:unnamed-chunk-9}Współczynnik koncentracji Giniego}
\centering
\begin{tabular}[t]{lr}
\toprule
Zmienna & Giniego\\
\midrule
\cellcolor{gray!10}{Price} & \cellcolor{gray!10}{0.708}\\
VolumeM & 0.989\\
\cellcolor{gray!10}{USAvgKWHPrice} & \cellcolor{gray!10}{0.055}\\
Difficulty & 0.742\\
\cellcolor{gray!10}{BTCCirculation} & \cellcolor{gray!10}{0.156}\\
\addlinespace
InflationRateUS & NA\\
\cellcolor{gray!10}{HasMuskMentioned} & \cellcolor{gray!10}{0.922}\\
GoogleTrends & 0.624\\
\bottomrule
\end{tabular}
\end{table}

\paragraph{Obserwacje}\label{obserwacje-8}

Analiza współczynnika Giniego dla badanych zmiennych ujawnia następujące
prawidłowości:

\begin{enumerate}
\def\labelenumi{\arabic{enumi}.}
\item
  \textbf{Cena Bitcoina (Price)} wykazuje bardzo wysoką koncentrację
  (G=0.708), co wskazuje na skrajną nierównomierność rozkładu wartości w
  czasie.
\item
  \textbf{Wolumen transakcji (VolumeM)} charakteryzuje się niemal
  maksymalną koncentracją (G=0.989), potwierdzając występowanie
  nielicznych okresów ekstremalnie wysokiej aktywności.
\item
  \textbf{Cena energii (USAvgKWHPrice)} prezentuje minimalną
  koncentrację (G=0.055), typową dla stabilnych wskaźników ekonomicznych
  o równomiernym rozkładzie.
\item
  \textbf{Trudność wydobycia (Difficulty)} ma wysoką wartość Giniego
  (0.742), odzwierciedlającą nieliniowy, skumulowany charakter wzrostu
  mocy obliczeniowej.
\item
  \textbf{Zmienne behawioralne}:

  \begin{itemize}
  \tightlist
  \item
    Tweety Muska (HasMuskMentioned) wykazują ekstremalną koncentrację
    (G=0.922)
  \item
    Wyszukiwania Google (GoogleTrends) mają wysoką koncentrację
    (G=0.624)
  \end{itemize}
\item
  \textbf{Podaż BTC (BTCCirculation)} charakteryzuje się umiarkowaną
  koncentracją (G=0.156), zgodną z przewidywalnym harmonogramem emisji.
\end{enumerate}

\paragraph{Wnioski}\label{wnioski-8}

\begin{enumerate}
\def\labelenumi{\arabic{enumi}.}
\tightlist
\item
  \textbf{Interpretacja ekonomiczna}: \vspace{0.2cm}

  \begin{itemize}
  \tightlist
  \item
    Wysokie wartości Giniego potwierdzają nierównomierny charakter
    rozwoju rynku krypto
  \item
    Niska koncentracja cen energii wzmacnia jej rolę jako stabilnego
    wskaźnika
  \item
    Ekstremalne wartości dla wolumenu i tweetów wskazują na ich
    impulsowy charakter
  \end{itemize}
\item
  \textbf{Implikacje metodologiczne}: \vspace{0.2cm}

  \begin{itemize}
  \tightlist
  \item
    Zmienne o G\textgreater0.6 wymagają specjalnych metod analitycznych
  \item
    Konieczność transformacji zmiennych o wysokiej koncentracji
  \item
    Wskazane jest modelowanie oddzielne dla okresów koncentracji
    wartości
  \end{itemize}
\item
  \textbf{Znaczenie dla zarządzania ryzykiem}: \vspace{0.2cm}

  \begin{itemize}
  \tightlist
  \item
    Wysoka koncentracja wskazuje na okresowe ``skupiska'' ekstremów
  \item
    Należy rozwijać strategie odporne na nierównomierność rozkładów
  \item
    Wskazane jest stosowanie analizy kwantylowej
  \end{itemize}
\end{enumerate}

\vspace{0.3cm}

\textit{Rekomendacje analityczne:}

\begin{itemize}
\item Dla zmiennych o G>0.8 zastosować analizę przedziałową
\item Przeprowadzić testy nierównomierności rozkładów
\item Uwzględnić efekty nieliniowe w modelowaniu
\end{itemize}

\pagebreak

\subsection{Analiza współczynników korelacji między
zmiennymi}\label{analiza-wspuxf3ux142czynnikuxf3w-korelacji-miux119dzy-zmiennymi}

\subsubsection{Macierz
współczynników}\label{macierz-wspuxf3ux142czynnikuxf3w}

\paragraph{Wyświetlenie tabeli}\label{wyux15bwietlenie-tabeli}

Poniższa tabela przedstawia macierz korelacji wszystkich zmiennych

\begin{table}[!h]
\centering
\caption{\label{tab:unnamed-chunk-12}Macierz korelacji zmiennych}
\centering
\resizebox{\ifdim\width>\linewidth\linewidth\else\width\fi}{!}{
\fontsize{9}{11}\selectfont
\begin{tabular}[t]{>{}lcccccccc}
\toprule
Zmienna & Price & VolumeM & USAvgKWHPrice & Difficulty & BTCCirculation & InflationRateUS & HasMuskMentioned & GoogleTrends\\
\midrule
\textbf{Price} & 1.0000 & 0.1084 & 0.8004 & 0.9139 & 0.6000 & 0.3860 & 0.2919 & 0.2256\\
\textbf{VolumeM} & 0.1084 & 1.0000 & 0.0489 & 0.0345 & 0.0864 & 0.2761 & 0.1097 & 0.0947\\
\textbf{USAvgKWHPrice} & 0.8004 & 0.0489 & 1.0000 & 0.9174 & 0.6394 & 0.4720 & 0.0745 & -0.0711\\
\textbf{Difficulty} & 0.9139 & 0.0345 & 0.9174 & 1.0000 & 0.5737 & 0.3305 & 0.1354 & -0.0453\\
\textbf{BTCCirculation} & 0.6000 & 0.0864 & 0.6394 & 0.5737 & 1.0000 & 0.3376 & 0.2341 & 0.3978\\
\textbf{InflationRateUS} & 0.3860 & 0.2761 & 0.4720 & 0.3305 & 0.3376 & 1.0000 & 0.1522 & 0.2216\\
\textbf{HasMuskMentioned} & 0.2919 & 0.1097 & 0.0745 & 0.1354 & 0.2341 & 0.1522 & 1.0000 & 0.4882\\
\textbf{GoogleTrends} & 0.2256 & 0.0947 & -0.0711 & -0.0453 & 0.3978 & 0.2216 & 0.4882 & 1.0000\\
\bottomrule
\end{tabular}}
\end{table}

\paragraph{Obserwacje}\label{obserwacje-9}

Analiza macierzy korelacji ujawnia następujące istotne zależności:

\begin{enumerate}
\def\labelenumi{\arabic{enumi}.}
\tightlist
\item
  \textbf{Silne korelacje dodatnie} (0.7-0.9):

  \begin{itemize}
  \tightlist
  \item
    Cena Bitcoina (Price) jest silnie skorelowana z:

    \begin{itemize}
    \tightlist
    \item
      Trudnością wydobycia (Difficulty, 0.914)
    \item
      Ceną energii w USA (USAvgKWHPrice, 0.800)
    \end{itemize}
  \item
    Cena energii (USAvgKWHPrice) i trudność wydobycia (Difficulty)
    wykazują niemal funkcjonalną zależność (0.917)
  \end{itemize}
\item
  \textbf{Umiarkowane korelacje} (0.4-0.7):

  \begin{itemize}
  \tightlist
  \item
    Podaż BTC (BTCCirculation) wykazuje umiarkowaną korelację z:

    \begin{itemize}
    \tightlist
    \item
      Ceną (0.600)
    \item
      Ceną energii (0.639)
    \end{itemize}
  \item
    Inflacja w USA (InflationRateUS) jest zauważalnie skorelowana z ceną
    energii (0.472)
  \end{itemize}
\item
  \textbf{Słabe korelacje} (\textless0.4):

  \begin{itemize}
  \tightlist
  \item
    Wszystkie zmienne behawioralne (HasMuskMentioned, GoogleTrends)
    wykazują słabe korelacje z ceną BTC
  \item
    Wolumen transakcji (VolumeM) jest praktycznie nieskorelowany z
    innymi zmiennymi
  \end{itemize}
\item
  \textbf{Ciekawe anomalie}:

  \begin{itemize}
  \tightlist
  \item
    GoogleTrends wykazuje ujemną korelację z ceną energii (-0.071) i
    trudnością wydobycia (-0.045)
  \item
    Korelacja między tweetami Muska a wyszukiwaniami Google (0.488) jest
    wyższa niż ich korelacje z ceną BTC
  \end{itemize}
\end{enumerate}

\paragraph{Wnioski}\label{wnioski-9}

\begin{enumerate}
\def\labelenumi{\arabic{enumi}.}
\tightlist
\item
  \textbf{Interpretacja ekonomiczna}: \vspace{0.2cm}

  \begin{itemize}
  \tightlist
  \item
    Silne korelacje cen energii i trudności wydobycia z ceną BTC
    sugerują dominację czynników fundamentalnych
  \item
    Słabe korelacje zmiennych behawioralnych wskazują na ich drugorzędne
    znaczenie w długim okresie
  \item
    Niska korelacja wolumenu z ceną podważa niektóre teorie techniczne
  \end{itemize}
\item
  \textbf{Implikacje dla modelowania}: \vspace{0.2cm}

  \begin{itemize}
  \tightlist
  \item
    Wysoka współliniowość cen energii i trudności wydobycia
    (VIF\textgreater5) wymaga selekcji zmiennych
  \item
    Zmienne behawioralne mogą wymagać nieliniowej specyfikacji w modelu
  \item
    Wskazane jest uwzględnienie opóźnień dla zmiennych makro
  \end{itemize}
\item
  \textbf{Rekomendacje analityczne}: \vspace{0.2cm}

  \begin{itemize}
  \item Zastosować analizę głównych składowych dla zmiennych silnie skorelowanych
  \item Przeprowadzić testy przyczynowości Grangera dla kluczowych zależności
  \item Rozważyć modele nieliniowe dla zmiennych behawioralnych
  \end{itemize}
\item
  \textbf{Ostrzeżenia}: \vspace{0.2cm}

  \begin{itemize}
  \tightlist
  \item
    Korelacja nie implikuje przyczynowości
  \item
    Wysokie wartości mogą wynikać ze wspólnego trendu czasowego
  \item
    Należy zbadać autokorelacje reszt
  \end{itemize}
\end{enumerate}

\vspace{0.3cm}

\textbf{Kluczowe zależności w postaci grafu:}

\begin{verbatim}
[USAvgKWHPrice] ←→ [Difficulty] → [Price]
      ↑               ↑             ↑
   [Inflation]     [BTC Supply]  [Behavioral]
\end{verbatim}

\subsection{Wykresy zależności między
zmiennymi}\label{wykresy-zaleux17cnoux15bci-miux119dzy-zmiennymi}

\subsubsection{Wyświetlenie
wykresów}\label{wyux15bwietlenie-wykresuxf3w}

\includegraphics{SPRAWOZDANIE_files/figure-latex/unnamed-chunk-13-1}
\includegraphics{SPRAWOZDANIE_files/figure-latex/unnamed-chunk-13-2}
\includegraphics{SPRAWOZDANIE_files/figure-latex/unnamed-chunk-13-3}
\includegraphics{SPRAWOZDANIE_files/figure-latex/unnamed-chunk-13-4}
\includegraphics{SPRAWOZDANIE_files/figure-latex/unnamed-chunk-13-5}
\includegraphics{SPRAWOZDANIE_files/figure-latex/unnamed-chunk-13-6}
\includegraphics{SPRAWOZDANIE_files/figure-latex/unnamed-chunk-13-7}

\subsubsection{Analiza wykresów}\label{analiza-wykresuxf3w}

\begin{itemize}
\item
  \textbf{Price vs VolumeM}: Widać pozytywną korelację -- wyższa cena
  powiązana jest z większym wolumenem. Relacja jednak nie jest idealnie
  liniowa.
\item
  \textbf{Price vs USAvgKWHPrice}: Niewielka dodatnia zależność; cena
  energii może mieć umiarkowany wpływ na cenę aktywa.
\item
  \textbf{Price vs Difficulty}: Bardzo silna pozytywna korelacja --
  większa trudność wydobycia powiązana z wyższą ceną.
\item
  \textbf{Price vs BTCCirculation}: Korelacja dodatnia, lecz słabsza;
  większa podaż w obiegu współwystępuje z wyższymi cenami.
\item
  \textbf{Price vs InflationRateUS}: Zależność niejednoznaczna; brak
  wyraźnego trendu.
\item
  \textbf{Price vs HasMuskMentioned}: Delikatna dodatnia zależność, ale
  z dużym rozrzutem -- możliwy wpływ medialny, ale słaby.
\item
  \textbf{Price vs GoogleTrends}: Widoczna dodatnia korelacja -- większe
  zainteresowanie w Google współwystępuje z wyższymi cenami.
\item
  \textbf{VolumeM vs USAvgKWHPrice}: Brak wyraźnej zależności -- wolumen
  nie reaguje istotnie na cenę energii.
\item
  \textbf{VolumeM vs Difficulty}: Silna dodatnia zależność -- większa
  trudność może zwiększać aktywność.
\item
  \textbf{VolumeM vs BTCCirculation}: Lekka zależność dodatnia --
  większa cyrkulacja łączy się z wyższym wolumenem.
\item
  \textbf{VolumeM vs InflationRateUS}: Brak widocznej korelacji.
\item
  \textbf{VolumeM vs HasMuskMentioned}: Zależność słaba, ale może być
  lekki trend wzrostowy.
\item
  \textbf{VolumeM vs GoogleTrends}: Wyraźna dodatnia zależność --
  większe zainteresowanie przekłada się na większy obrót.
\item
  \textbf{USAvgKWHPrice vs Difficulty}: Zależność bliska zeru -- brak
  widocznego wpływu ceny energii na trudność.
\item
  \textbf{USAvgKWHPrice vs BTCCirculation}: Brak wyraźnej zależności.
\item
  \textbf{USAvgKWHPrice vs InflationRateUS}: Słaba dodatnia korelacja --
  może istnieć pośredni wpływ inflacji na koszty energii.
\item
  \textbf{USAvgKWHPrice vs HasMuskMentioned}: Bez zauważalnego trendu.
\item
  \textbf{USAvgKWHPrice vs GoogleTrends}: Również brak wyraźnego wzoru.
\item
  \textbf{Difficulty vs BTCCirculation}: Silna dodatnia zależność -- im
  więcej BTC w obiegu, tym trudniej je wydobyć.
\item
  \textbf{Difficulty vs InflationRateUS}: Brak korelacji -- czynniki
  wewnętrzne sieci dominują.
\item
  \textbf{Difficulty vs HasMuskMentioned}: Nie obserwuje się wyraźnej
  zależności.
\item
  \textbf{Difficulty vs GoogleTrends}: Widoczna korelacja -- trudność
  wydobycia rośnie z zainteresowaniem publicznym.
\item
  \textbf{BTCCirculation vs InflationRateUS}: Brak zależności.
\item
  \textbf{BTCCirculation vs HasMuskMentioned}: Korelacja niewielka,
  raczej przypadkowa.
\item
  \textbf{BTCCirculation vs GoogleTrends}: Lekka dodatnia zależność --
  większe zainteresowanie może prowadzić do większej cyrkulacji.
\item
  \textbf{InflationRateUS vs HasMuskMentioned}: Brak widocznego związku.
\item
  \textbf{InflationRateUS vs GoogleTrends}: Nie obserwuje się korelacji.
\item
  \textbf{HasMuskMentioned vs GoogleTrends}: Zależność dodatnia --
  więcej wzmiankowań Elona Muska często pokrywa się ze wzrostem
  wyszukiwań.
\end{itemize}

\pagebreak

\section{Przygotowanie danych}\label{przygotowanie-danych}

\subsection{Stacjonarność}\label{stacjonarnoux15bux107}

\subsubsection{Badanie stacjonarności
zmiennych}\label{badanie-stacjonarnoux15bci-zmiennych}

\begingroup\fontsize{10}{12}\selectfont

\begin{longtable}[t]{llcccl}
\caption{\label{tab:unnamed-chunk-14}Wyniki testów stacjonarności ADF i KPSS dla zmiennych modelu}\\
\toprule
 & Zmienna & Test & Statystyka testu & p-value & Wniosek\\
\midrule
\addlinespace[0.3em]
\multicolumn{6}{l}{\textbf{Price}}\\
\hspace{1em}Dickey-Fuller...1 & Price & ADF & -0,105 & 0.9900 & Niestacjonarna\\
\hspace{1em}KPSS Level...2 & Price & KPSS & 2,480 & 0.0100 & Niestacjonarna\\
\addlinespace[0.3em]
\multicolumn{6}{l}{\textbf{VolumeM}}\\
\hspace{1em}Dickey-Fuller...3 & VolumeM & ADF & -5,177 & 0.0100 & Stacjonarna\\
\hspace{1em}KPSS Level...4 & VolumeM & KPSS & 0,192 & 0.1000 & Stacjonarna\\
\addlinespace[0.3em]
\multicolumn{6}{l}{\textbf{USAvgKWHPrice}}\\
\hspace{1em}Dickey-Fuller...5 & USAvgKWHPrice & ADF & -0,891 & 0.9517 & Niestacjonarna\\
\hspace{1em}KPSS Level...6 & USAvgKWHPrice & KPSS & 2,505 & 0.0100 & Niestacjonarna\\
\addlinespace[0.3em]
\multicolumn{6}{l}{\textbf{Difficulty}}\\
\hspace{1em}Dickey-Fuller...7 & Difficulty & ADF & 4,482 & 0.9900 & Niestacjonarna\\
\hspace{1em}KPSS Level...8 & Difficulty & KPSS & 2,441 & 0.0100 & Niestacjonarna\\
\addlinespace[0.3em]
\multicolumn{6}{l}{\textbf{BTCCirculation}}\\
\hspace{1em}Dickey-Fuller...9 & BTCCirculation & ADF & -4,497 & 0.0100 & Stacjonarna\\
\hspace{1em}KPSS Level...10 & BTCCirculation & KPSS & 3,273 & 0.0100 & Niestacjonarna\\
\addlinespace[0.3em]
\multicolumn{6}{l}{\textbf{InflationRateUS}}\\
\hspace{1em}Dickey-Fuller...11 & InflationRateUS & ADF & -2,539 & 0.3511 & Niestacjonarna\\
\hspace{1em}KPSS Level...12 & InflationRateUS & KPSS & 1,159 & 0.0100 & Niestacjonarna\\
\addlinespace[0.3em]
\multicolumn{6}{l}{\textbf{HasMuskMentioned}}\\
\hspace{1em}Dickey-Fuller...13 & HasMuskMentioned & ADF & -4,260 & 0.0100 & Stacjonarna\\
\hspace{1em}KPSS Level...14 & HasMuskMentioned & KPSS & 0,707 & 0.0129 & Niestacjonarna\\
\addlinespace[0.3em]
\multicolumn{6}{l}{\textbf{GoogleTrends}}\\
\hspace{1em}Dickey-Fuller...15 & GoogleTrends & ADF & -2,611 & 0.3208 & Niestacjonarna\\
\hspace{1em}KPSS Level...16 & GoogleTrends & KPSS & 0,844 & 0.0100 & Niestacjonarna\\
\bottomrule
\end{longtable}
\endgroup{}

\subsubsection{Eliminacja
niestacjonarności}\label{eliminacja-niestacjonarnoux15bci}

dużo rzeczy

\subsection{Logarytmowanie}\label{logarytmowanie}

\subsubsection{Wybór zmiennej do
zlogarytmowania}\label{wybuxf3r-zmiennej-do-zlogarytmowania}

\subsubsection{Dodanie kolumny zmiennej
logarytmowanej}\label{dodanie-kolumny-zmiennej-logarytmowanej}

\subsubsection{Weryfikacja}\label{weryfikacja}

\subsection{Eliminacja wartości
odstających}\label{eliminacja-wartoux15bci-odstajux105cych}

bla bla bla na.rm i rm.outliers i już

\section{Zastosowanie metod doboru zmiennych do
modelu}\label{zastosowanie-metod-doboru-zmiennych-do-modelu}

\subsection{Metoda Hellwiga}\label{metoda-hellwiga}

bim bam bom dużo liczenia i podpunktów

\subsection{Metoda krokowa wsteczna}\label{metoda-krokowa-wsteczna}

bim bam bom dużo liczenia i podpunktów znowuuuu

\section{Budowa modelu}\label{budowa-modelu}

\subsection{Podział danych na część uczącą i
testową}\label{podziaux142-danych-na-czux119ux15bux107-uczux105cux105-i-testowux105}

cośtam trzeba o tym powiedzieć i napisać

\subsection{Budowa modelu regresji liniowej na danych z części
uczącej}\label{budowa-modelu-regresji-liniowej-na-danych-z-czux119ux15bci-uczux105cej}

podać kod czy coś

\section{Weryfikacja poprawności
modelu}\label{weryfikacja-poprawnoux15bci-modelu}

\subsection{Badanie normalności rozkładu
reszt}\label{badanie-normalnoux15bci-rozkux142adu-reszt}

dużo danych więc okej, dodać testy i nie musi przechodzic przez testy
żeby było ściśle super

\subsection{Testowanie autokorelacji}\label{testowanie-autokorelacji}

\subsection{Badanie
heteroskedastyczności}\label{badanie-heteroskedastycznoux15bci}

\subsubsection{Test Goldfelda-Quandta}\label{test-goldfelda-quandta}

\paragraph{Obliczenie}\label{obliczenie-6}

bum

\paragraph{Obserwacje}\label{obserwacje-10}

bum

\paragraph{Wnioski}\label{wnioski-10}

bum

\subsubsection{Test Breuscha-Pagana}\label{test-breuscha-pagana}

\paragraph{Obliczenie}\label{obliczenie-7}

bum

\paragraph{Obserwacje}\label{obserwacje-11}

bum

\paragraph{Wnioski}\label{wnioski-11}

bum

\subsubsection{Test White'a}\label{test-whitea}

\paragraph{Obliczenie}\label{obliczenie-8}

bum

\paragraph{Obserwacje}\label{obserwacje-12}

bum

\paragraph{Wnioski}\label{wnioski-12}

bum

\subsection{Testowanie współliniowości
(VIF)}\label{testowanie-wspuxf3ux142liniowoux15bci-vif}

\subsection{Testowanie stabilności parametrów
modelu}\label{testowanie-stabilnoux15bci-parametruxf3w-modelu}

\subsubsection{Test Chowa}\label{test-chowa}

\paragraph{Obliczenie}\label{obliczenie-9}

bum

\paragraph{Obserwacje}\label{obserwacje-13}

bum

\paragraph{Wnioski}\label{wnioski-13}

bum

\subsection{Testowanie stabilności postaci analitycznej
modelu}\label{testowanie-stabilnoux15bci-postaci-analitycznej-modelu}

\subsubsection{Test Ramsey'a RESET}\label{test-ramseya-reset}

\paragraph{Obliczenie}\label{obliczenie-10}

bum

\paragraph{Obserwacje}\label{obserwacje-14}

bum

\paragraph{Wnioski}\label{wnioski-14}

bum

\subsubsection{Test liczby serii}\label{test-liczby-serii}

\paragraph{Obliczenie}\label{obliczenie-11}

bum

\paragraph{Obserwacje}\label{obserwacje-15}

bum

\paragraph{Wnioski}\label{wnioski-15}

bum

\subsection{Badanie efektu katalizy}\label{badanie-efektu-katalizy}

bum bum

\subsection{Badanie koincydencji}\label{badanie-koincydencji}

\section{Prezentacja ostatecznej postaci modelu i jego
ocena}\label{prezentacja-ostatecznej-postaci-modelu-i-jego-ocena}

\subsection{Ocena istotności
zmiennych}\label{ocena-istotnoux15bci-zmiennych}

\subsubsection{Test t-Studenta}\label{test-t-studenta}

\paragraph{Obliczenie}\label{obliczenie-12}

bum

\paragraph{Obserwacje}\label{obserwacje-16}

bum

\paragraph{Wnioski}\label{wnioski-16}

bum

\subsubsection{Test Walda}\label{test-walda}

\paragraph{Obliczenie}\label{obliczenie-13}

bum

\paragraph{Obserwacje}\label{obserwacje-17}

bum

\paragraph{Wnioski}\label{wnioski-17}

bum

\subsection{Ocena współczynnika
determinacji}\label{ocena-wspuxf3ux142czynnika-determinacji}

bum bum

\subsection{Interpretacja parametrów
modelu}\label{interpretacja-parametruxf3w-modelu}

tu będzie dużo pisania i sporo podpunktów

\section{Prognoza}\label{prognoza}

\subsection{Przeprowadzenie prognozy
punktowej}\label{przeprowadzenie-prognozy-punktowej}

\subsection{Sprawdzenie błędu prognozy EX
POST}\label{sprawdzenie-bux142ux119du-prognozy-ex-post}

\end{document}
